% Provenance: zvf-program/theory/zvf_theory.tex (T1-T3, post-correction),
% analyze_t1_ci.py / analyze_t2_floor.py / analyze_t3_gstar.py artifacts.
\chapter{Estimation Theory and Reliability Budgets}
\label{ch:theory}

This chapter presents the theory in its \emph{corrected} form. Two errors in
the programme's earlier statements --- both found by external adversarial
review on 2026-07-11 and verified against the source --- are reported openly
in Section~\ref{sec:corrections}, because the corrections are themselves
informative.

\section{T1: \ZVF{} as a calibrated estimator}
$\ZVF_t$ over $m$ groups is an unbiased binomial-proportion estimator of the
population uniformity rate $\theta_G=\mathbb{E}_\phi[h_G(p)]$, with
$h_G(p)=p^G+(1-p)^G$ and $\phi$ the prompt-difficulty prior. Empirical
coverage on 512-prompt pools: the Wald interval covers 0.92--0.97 at
$m\ge64$ (marginal at $m{=}32$); the \textbf{Wilson interval covers
0.95--0.98 in every tested setting} and is adopted operationally.
Under difficulty-\emph{sorted} (curriculum) batching, global coverage
collapses to 0.08--0.40 --- but the interval remains calibrated (0.944) for
the \emph{local}, curriculum-stage \ZVF{}, and stratified batch composition
restores conservative global validity (0.996). The across-group-independence
threat is thus an estimand-labelling requirement, not an invalidation: report
\emph{which} \ZVF{} the interval covers.

\section{T2: the waiting-time reliability budget}
With i.i.d.\ groups, each uninformative with probability \ZVF{}, the number of
groups to the first informative one is geometric, giving
\[
N(\ZVF) \;=\; G\left\lceil \frac{\ln \delta}{\ln \ZVF}\right\rceil
\]
as the $(1{-}\delta)$-quantile of rollouts-to-first-informative-group.
\textbf{This is a reliability budget, not an impossibility bound}: an
informative group arrives within the first $G$ rollouts with probability
$1-\ZVF$, which can be large; a budget below $N$ does not make a gradient
improbable, it makes the \emph{guarantee} of one unavailable at confidence
$1-\delta$. Empirically the geometric model is essentially exact: with
$G{=}8$, $\delta{=}0.1$, the model quantile matched the observed quantile at
ratio $1.00$ in all six difficulty strata; in the hardest stratum
($\hat p\approx0.01$, $\ZVF{=}0.886$) the budget is 160 rollouts. Composed
with T1, the one-sided lower confidence bound $\underline{\ZVF}_t$ yields a
per-step, observation-driven budget $N(\underline{\ZVF}_t)$ --- conservative,
since $N$ is increasing and $N(\underline{\ZVF}_t)\le N(\ZVF_t)$.

\section{T3: the optimal group size is universal (a negative result)}
The signal-per-rollout objective
$J(G)=\tfrac1G\,\mathbb{E}_{p\sim\phi}[(1-h_G(p))\,p(1-p)]$ was intended to
yield a prior-dependent optimal group size. It does not. The identities
$1-p^2-(1-p)^2=2p(1-p)$ and $1-p^3-(1-p)^3=3p(1-p)$ hold pointwise, so
\[
J(2)\;=\;J(3)\;=\;\mathbb{E}_{p\sim\phi}\big[(p(1-p))^2\big]
\quad\text{for every prior } \phi,
\]
and $J(G)\le J(2)$ for $G\ge4$. Hence $G^\star\in\{2,3\}$ \emph{universally}:
the objective's argmax carries no information about the data, the previously
published ``empirical argmax 2--3 agrees with theory'' validation was
algebraically guaranteed rather than a test, and any data-adaptive
group-size policy must come from a richer objective (variance-, Fisher-, or
wall-clock-aware). We report this as a negative result of independent value:
it explains why per-rollout accounting always favours tiny groups, and it
sharpens Chapter~6's message that the real trade-off is temporal (early
step-count versus endgame signal), not a static optimum.

\section{Corrections adopted from external review}
\label{sec:corrections}
(1) An earlier statement of T2 claimed ``no nonzero gradient step is possible
with fewer than $N$ rollouts except with probability $<\delta$'' --- a
quantifier error confusing a quantile with a minimum; the validation script
additionally tested the inverted inequality. Both are corrected, the analysis
was re-run (fits at ratio 1.00 unchanged), and the corollary's
confidence-direction commentary was rewritten (a lower confidence bound on
\ZVF{} gives a \emph{smaller}, conservative budget). (2) The theory draft
briefly defined $\GU_t:=\ZVF_t$, contradicting the programme-wide
$\GU=1-\ZVF$; fixed. The proof-gap registers of the source paper (mixture
versus fixed-$p$ geometric modelling; improvement versus nonzero-gradient;
the global $G\ge4$ inequality) remain open and are inherited here as stated
limitations.
